\subsection{普遍正核}

在本章中，X 是一个分离拓扑空间。

定理 1. --- 设 \(f \in \mathcal{C}(X \times X)\)。下列条件等价：

\begin{enumerate}
\def\labelenumi{(\roman{enumi})}
\tightlist
\item
  对 X 的每个紧子集 Y 及 Y 上的每个正测度 \(\mu\)，\(\mathrm{L}^{2}(\mathrm{Y},\mu)\) 的由核 \(f|(\mathrm{Y}\times\mathrm{Y})\)（III, 第 29 页, 定义 1）定义的自同态是正的，换句话说，我们有
\end{enumerate}

\[
\int_ {\mathrm{Y} \times \mathrm{Y}} \overline {{h (x)}} h (y) f (x, y) d (\mu \otimes \mu) (x, y) \geqslant 0
\]

对所有 \(h \in \mathcal{L}^{2}(\mathrm{Y}, \mu)\) 成立；

\begin{enumerate}
\def\labelenumi{(\roman{enumi})}
\setcounter{enumi}{1}
\tightlist
\item
  对每个整数 \(n \in N\)、X 中每个族 \((x_{i})_{0 \leqslant i \leqslant n}\) 以及复数的任意族 \((t_{i})_{0 \leqslant i \leqslant n}\)，我们有
\end{enumerate}

\[
\sum_ {i = 0} ^ {n} \sum_ {j = 0} ^ {n} \bar {t} _ {i} t _ {j} f (x _ {i}, x _ {j}) \geqslant 0;
\]

\begin{enumerate}
\def\labelenumi{(\roman{enumi})}
\setcounter{enumi}{2}
\item
  存在一个复希尔伯特空间 E 和一个具有完全系像的连续映射 \(g \colon X \to E\)，使得 \(f(x, y) = \langle g(x) | g(y) \rangle\) 对所有 \(x\) 与 \(y\)（\(X\) 中的）成立；
\item
  存在一个复希尔伯特空间 E 和一个连续映射 \(g: X \to E\)，使得对所有 X 中的 x 与 y 有 \(f(x, y) = \langle g(x) | g(y) \rangle\)。
\end{enumerate}

(iii) 蕴含 (iv) 是显然的，而通过考虑离散测度 \(\mu\)，即 \(\{0,\ldots,n\}\) 上的计数测度经映射 \(i\mapsto x_{i}\) 的像，以及满足下式的函数 h，可知 (i) 蕴含 (ii)

\[
h(x) = \sum_{\substack{0\leqslant i\leqslant n\\ x_{i} = x}}t_{i}.
\]

我们来证明 (iv) 蕴含 (i)。设存在一个复希尔伯特空间 E 和一个连续映射 \(g \colon \mathrm{X} \to \mathrm{E}\)，使得 \(f(x, y) = \langle g(x) \mid g(y) \rangle\) 对所有 \((x, y) \in \mathrm{X} \times \mathrm{X}\)。设 Y 是 X 的一个紧子集，\(\mu\) 是 Y 上的一个正测度，且 \(h \in \mathcal{L}^{2}(Y, \mu)\)。我们有

\[
\begin{array}{r l} & {\int_ {\mathrm{Y} \times \mathrm{Y}} \overline {{h (x)}} h (y) f (x, y) d (\mu \otimes \mu) (x, y)} \\ & {\qquad = \int_ {\mathrm{Y} \times \mathrm{Y}} \overline {{h (x)}} h (y) \langle g (x) | g (y) \rangle d (\mu \otimes \mu) (x, y)} \\ & {\qquad = \left\langle \int_ {\mathrm{Y}} h (x) g (x) d \mu (x) \Big | \int_ {\mathrm{Y}} h (y) g (y) d \mu (y) \right\rangle \geqslant 0} \end{array}
\]

（INT, V, 第 97 页, § 8, 第 4 目, 命题 9）。

最后我们来证明 (ii) 蕴含 (iii)。设 \(\widetilde{\mathrm{E}}\) 是 X 上具有限支撑的复测度组成的空间。对 \(\mu_{1}\) 与 \(\mu_{2}\)（\(\widetilde{\mathrm{E}}\) 中的），我们令

\[
\langle \mu_ {1} \mid \mu_ {2} \rangle = \int_ {\mathrm{X} \times \mathrm{X}} f (x, y) (\overline {{\mu}} _ {1} \otimes \mu_ {2}) (x, y).
\]

这样在 \(\widetilde{\mathrm{E}}\) 上定义的半双线性形式是一个正埃尔米特形式。事实上，设 \(\mu \in \widetilde{\mathrm{E}}\)；存在 X 中的有限族 \((x_{i})_{0\leqslant i\leqslant n}\) 及复数 \((t_i)_{0\leqslant i\leqslant n}\)，使得 \(\mu = \sum_{i=0}^{n} t_i \varepsilon_{x_i}\)。于是我们有

\[
\langle \mu | \mu \rangle = \sum_ {i = 0} ^ {n} \sum_ {j = 0} ^ {n} \bar {t} _ {i} t _ {j} f (x _ {i}, x _ {j}) \geqslant 0
\]

（由假设）。定义 \(\widetilde{g}\colon\mathbf{X}\to\widetilde{\mathbf{E}}\) 为 \(\widetilde{g}(x)=\varepsilon_{x}\)。映射 \(\widetilde{g}\) 的像生成 \(\widetilde{\mathbf{E}}\)。此外，对每个 \((x,y)\in\mathbf{X}\times\mathbf{X}\)，一方面有 \(f(x,y)=\langle\widetilde{g}(x)|\widetilde{g}(y)\rangle\)，另一方面有

\[
\| \widetilde {g} (x) - \widetilde {g} (y) \| ^ {2} = f (x, x) + f (y, y) - f (x, y) - f (y, x),
\]

这蕴含 \(\widetilde{g}\) 连续，因为 \(f\) 连续。

设 E 是 \(\widetilde{E}\) 的分离完备希尔伯特空间（EVT, V, 第 8 页, 命题 4 的推论），\(g\colon X \to E\) 是 \(\widetilde{g}\) 与典范映射 \(\widetilde{E} \to E\) 的复合。则映射 g 连续，它的像在 E 中是完全系，且 \(f(x,y) = \langle g(x) | g(y) \rangle\) 对所有 \((x,y) \in X \times X\)。

用来证明 (ii) 蕴含 (iii) 的方法称为盖尔范德–奈马克–塞加尔构造（GNS 构造）。

定义 1. --- 若函数 \(f \in \mathcal{C}(\mathrm{X} \times \mathrm{X})\) 满足定理 1 的等价条件，则称它是 X 上的一个普遍正核。

若 f 是 X 上的一个普遍正核，则满足出处同上条件 (iv) 的偶 \((E,g)\) 称为 f 的一个希尔伯特实现；若条件 (iii) 得到满足，则称它是一个循环希尔伯特实现。

用 \(\mathrm{Noy}_{+}(\mathrm{X})\) 表示 X 上普遍正核的集合。

设 X' 是一个分离拓扑空间，\(h: X \to X'\) 是一个连续映射。映射 \(f \mapsto f \circ (h, h)\)（\(\mathcal{C}(X' \times X')\) 到 \(\mathcal{C}(X \times X)\) 中的）通过过渡到子空间诱导出映射 \(\mathrm{Noy}_{+}(X') \to \mathrm{Noy}_{+}(X)\)。

命题 1. --- 设 f 是 X 上的一个普遍正核。设 \((\mathrm{E}_{1}, g_{1})\) 与 \((\mathrm{E}_{2}, g_{2})\) 是 f 的希尔伯特实现。假设 \((\mathrm{E}_{1}, g_{1})\) 是一个循环希尔伯特实现。则存在唯一的连续线性映射 \(u: E_{1} \to E_{2}\) 使得 \(g_{2} = u \circ g_{1}\)。这个映射是等距的。若 \((\mathrm{E}_{2}, g_{2})\) 也是循环的，则 u 是一个同构。

\(u\) 的唯一性由 \(g_{1}\) 的像在 \(\mathrm{E}_{1}\) 中是完全系这一事实得出。设 \(\mathrm{F} = \mathbf{C}^{(\mathrm{X})}\)，\((e_x)_{x\in \mathrm{X}}\) 是 \(\mathrm{F}\) 的典范基。对 \(j = 1\) 与 \(j = 2\)，用 \(u_{j}\) 表示从 \(\mathrm{F}\) 到 \(\mathrm{E}_{j}\) 中的由 \(u_{j}(e_{x}) = g_{j}(x)\) 确定的线性映射，并用 \(\mathrm{F}_{j}\) 表示它的像；空间 \(\mathrm{F}_{1}\) 在 \(\mathrm{E}_{1}\) 中稠密。设 \(t = \sum t_{x}e_{x}\) 是 \(\mathrm{F}\) 的一个元素。我们有

\[
\begin{array}{c} \| u _ {1} (t) \| ^ {2} = \sum_ {x, y} \bar {t} _ {x} t _ {y} \langle g _ {1} (x)   |   g _ {1} (y) \rangle = \sum_ {x, y} \bar {t} _ {x} t _ {y} f (x, y) \\ = \sum_ {x, y} \bar {t} _ {x} t _ {y} \langle g _ {2} (x)   |   g _ {2} (y) \rangle = \| u _ {2} (t) \| ^ {2}. \end{array}
\]

因此，存在从 \(F_{1}\) 到 \(F_{2}\) 中的线性等距映射 v 使得 \(u_{2}=v\circ u_{1}\)，特别地，\(g_{2}(x)=v(g_{1}(x))\) 对所有 \(x\in X\)。由于 \(g_{1}\) 的像在 \(E_{1}\) 中是完全系，这个映射延拓为从 \(E_{1}\) 到 \(E_{2}\) 中的等距线性映射 u，使得 \(g_{2}=u\circ g_{1}\)。

由 I 第 107 页的引理 8，\(u\) 的像在 \(\mathrm{E}_2\) 中是闭的。若 \((\mathrm{E}_2,g_2)\) 是一个循环希尔伯特实现，则 \(u\) 的像还在 \(\mathrm{E}_2\) 中稠密，因而 \(u\) 是一个同构。

命题 2. --- 集合 Noy \(_{+}\) (X) 是空间 \(\mathcal{C}(\text{X} \times \text{X})\) 中的一个自伴锥；它在乘法下稳定，且当给 \(\mathcal{C}(\text{X} \times \text{X})\) 赋以逐点收敛拓扑时它是闭的。

若 \(f \in \mathrm{Noy}_{+}(\mathrm{X})\)，则 \(tf \in \mathrm{Noy}_{+}(\mathrm{X})\) 对每个实数 \(t \geqslant 0\)，且 \(\overline{f} \in \mathrm{Noy}_{+}(\mathrm{X})\)，这是初等的。

若 \((\mathrm{E}_{1}, g_{1})\) 与 \((\mathrm{E}_{2}, g_{2})\) 分别是 X 上普遍正核 \(f_{1}\) 与 \(f_{2}\) 的希尔伯特实现，则偶 \((\mathrm{E}_{1} \oplus \mathrm{E}_{2}, g_{1} + g_{2})\)（相应地，偶 \((\mathrm{E}_{1} \widehat{\otimes}_{2} \mathrm{E}_{2}, g_{1} \otimes g_{2})\)）是 \(f_{1} + f_{2}\)（相应地 \(f_{1}f_{2}\)）的一个希尔伯特实现；因而它们是普遍正核。

\(\mathrm{Noy}_{+}(\mathrm{X})\) 的刻画 (ii)（定理 1）蕴含这个集合在赋以逐点收敛拓扑的 \(\mathcal{C}(\mathrm{X}\times\mathrm{X})\) 中是闭的。

